\section{Comparing finite integer averages with independent residues}
\label{q:finite-law-section}

The auxiliary law makes the prime residue coordinates independent,
but the prime sets grow with $X$.  Averaging over their full joint
period would therefore give no useful error at the required scale.
We instead compare the particular residue tests used in the proof.
The comparison rests on bounded independence for Boolean circuits;
we give the reduction from integer residues, including the small
correction needed to make all sufficiently small bit marginals exactly uniform.

For a consecutive interval $I$ of integers, write $\E_I$ for the
average obtained by choosing the origin uniformly in $I$.
The symbol $\E$ continues to denote the independent residue law,
and $R_s$ denotes its coordinate modulo $s$.
A translated residue test at offset $a$ means
$R_s\equiv-a\pmod s$, using the same coordinate $R_s$ at every
translate.  Different sites do not receive independent new coordinates.
A Boolean circuit here uses unbounded-fan-in AND and OR gates and
negations.  Its depth counts AND and OR levels; for its size we count
both gates and input occurrences.  An event is represented by a circuit
whose input bits are residue equalities.

\begin{lemma}[Finite residue comparison]\label{q:finite-law}
There is an absolute constant $C_1$ with the following property.
Choose $A>C_1+20$.  Let $l\ge1$ be fixed, and let $\mathcal R$ be any
set of primes at most $e^L$, none dividing $l$.  On the coordinates
modulo $p\in\mathcal R$ and modulo $l$, let $\mu$ be the product of
the uniform laws; the modulus-one coordinate can be omitted.
Suppose that an event $\mathcal E$ is represented by a circuit of
depth at most $20$ and size at most $\exp(L^6)$ in these residue
equalities.  For every consecutive integer interval $I$ with
$|I|\ge\exp(L^A/2)$,
\begin{equation}\label{q:finite-law-bound}
 \bigl|\E_I\1_{\mathcal E}-\mu(\mathcal E)\bigr|
       \le e^{-L^{10}}
\end{equation}
for sufficiently large $L$ depending on $l$.
The bound is uniform in the position of $I$ and in all offsets used
by the tests.  The modulus $l$ need not be squarefree.
\end{lemma}

\begin{proof}
We use Braverman's bounded-independence theorem
\cite[Corollary~2]{Braverman10}: for a fixed depth $d_0$, every
$t$-wise uniform distribution on Boolean bits fools circuits of
size $m$ to error $\varepsilon$ when
$t\ge(\log(m/\varepsilon))^{C(d_0)}$, with a suitable constant
depending only on the depth.  Being $t$-wise uniform means that every
set of at most $t$ bits has exactly the uniform joint distribution.
We first encode the residue tests by bits, then correct the small
failure of exact independence in the integer average.

Under the integer-residue law, $R_s$ denotes the residue of the
uniformly chosen origin modulo $s$; under $\mu$ it is the independent
uniform coordinate. In both cases represent $R_s$ by an integer in
$\{0,\ldots,s-1\}$. Put $B=\lceil L^{15}\rceil$. Independently at each modulus $s$,
adjoin a uniform random variable $U_s\in[0,1)$, independent also of
the origin, and take the first $B$ binary digits of
\[
                   Z_s=\frac{R_s+U_s}{s}.
\]
Under $\mu$, all these bits are independent and unbiased.  Decode
$R_s$ from its bits whenever their dyadic interval lies inside one
interval $[r/s,(r+1)/s)$; on ambiguous intervals make any fixed choice.
Conditional on any value of $R_s$, the probability of a decoding
error is at most $2s2^{-B}$, since only the two boundary cells can
cause an error.  This bound holds under either origin law.  Therefore
the probability of any error is at most
\begin{equation}\label{q:decoding-error}
 2^{1-B}\left(l+\sum_{p\in\mathcal R}p\right)
                 \le e^{3L}2^{-B}
\end{equation}
for large $L$.  Correct decoding of one coordinate makes all its
translated equality tests correct simultaneously.

A decoded equality is a Boolean function of $B$ bits.  Its
truth-table disjunctive normal form has at most $2^B$ conjunctions,
each containing at most $B$ literals.  Substituting these forms into
the original circuit gives depth at most $22$ and size at most
$\exp(L^{17})$ for large $L$.

Fix any set of at most $t$ encoded bits.  It involves at most $t$
residue coordinates, whose moduli are pairwise coprime.  Their
product $D$ satisfies $D\le l e^{tL}$.  If $N=|I|$, each residue
modulo $D$ occurs either $\lfloor N/D\rfloor$ or $\lceil N/D\rceil$
times in $I$.  Thus its distribution has total variation distance
at most $D/N$ from uniform.  Here total variation is one half of
the $\ell^1$ distance between probability masses.  Adding independent
jitters and projecting to the selected bits cannot increase this
distance.  Their marginal consequently differs from uniform by at
most
\begin{equation}\label{q:crt-bits}
 \Delta\le \exp(2tL-L^A/2)
\end{equation}
for sufficiently large $L$.  The only use of $l$ is as one modulus
coprime to the selected primes, so repeated prime factors in $l$
cause no change.

There are at most $n_{\mathrm{bit}}\le e^{2L}$ bits in all.
Let $f$ be the density of their integer-origin law relative to
uniform measure on this finite Boolean cube.  For a set $S$ of bit
indices, let $\chi_S$ be the product of the signs $(-1)^{x_i}$ over
$i\in S$, and write $\widehat f(S)=\E_{\mathrm{unif}}f\chi_S$.
By \eqref{q:crt-bits},
\[
                  |\widehat f(S)|\le2\Delta
       \qquad(1\le |S|\le t).
\]
Set $t=\lceil L^{C_1}\rceil$.  The total absolute value of these
coefficients is bounded by
\begin{equation}\label{q:small-fourier-mass}
 q:=\sum_{1\le|S|\le t}|\widehat f(S)|
 \le 2\Delta(t+1)n_{\mathrm{bit}}^t
 \le e^{-L^{11}},
\end{equation}
provided $A>C_1+20$ and $L$ is large.

The estimate \eqref{q:small-fourier-mass} allows a correction that
preserves nonnegativity, even where $f$ vanishes.  This is a direct
Fourier form of the almost-to-exact independence argument of Alon,
Goldreich, and Mansour \cite[Theorem~2.1 and Section~3.1]{AGM03}.
We give the calculation with the error budget needed here.
Put $a=e^{-L^{11}}$
and define
\[
 H=\sum_{1\le|S|\le t}\widehat f(S)\chi_S,
 \qquad g=\frac{f-H+a}{1+a}.
\]
Since $|H|\le q\le a$, the function $g$ is nonnegative.  Its integral
is one, and every nonconstant Fourier coefficient of order at most
$t$ vanishes.  Fourier inversion on each marginal shows that the
law with density $g$ is exactly $t$-wise uniform.  Moreover,
\begin{equation}\label{q:fourier-correction-tv}
 d_{\mathrm{TV}}(f,g)
 \le\frac{q+a\E_{\mathrm{unif}}|1-f|}{2(1+a)}
 \le\frac{3a}{2}.
\end{equation}

Choose the absolute constant $C_1$ large enough for Braverman's
theorem at depth $22$, size $e^{L^{17}}$, and error $e^{-L^{11}}$.
Then choose $A>C_1+20$. Applying Braverman's theorem to $g$ gives
error at most $a=e^{-L^{11}}$. Under either the integer-residue or
product-residue law, the original and encoded tests can disagree only
on the decoding event. The total comparison error is therefore at most
\[
 2e^{3L}2^{-B}+d_{\mathrm{TV}}(f,g)+a
 \le 2e^{3L}2^{-B}+\tfrac52e^{-L^{11}}
 \le e^{-L^{10}}
\]
for sufficiently large $L$. This proves \eqref{q:finite-law-bound}.
All choices are absolute; the threshold may depend on the fixed modulus
$l$.
\end{proof}

\begin{corollary}[Scalar expansions]\label{q:finite-expansion}
Under the hypotheses of \Cref{q:finite-law}, write $\E$ for
expectation under $\mu$.  Suppose
$\Phi=\sum_\alpha c_\alpha\1_{\mathcal E_\alpha}$, where every event
has the stated circuit bounds and
$\sum_\alpha|c_\alpha|\le\exp(C L^5)$ for a fixed constant $C$.
Then, for sufficiently large $L$,
\begin{equation}\label{q:scalar-comparison}
                 |\E_I\Phi-\E\Phi|\le e^{-L^9}.
\end{equation}
This also applies to the sum of errors for a family of such expansions
when the same bound holds for their total absolute coefficient sum.
\end{corollary}

\begin{proof}
Apply \eqref{q:finite-law-bound} term by term and sum
$|c_\alpha|e^{-L^{10}}$.  This argument preserves the signs of the
coefficients until after taking each expectation.
\end{proof}

\subsection{Low-degree residue states and truncated weights}

To apply \Cref{q:finite-expansion} to weights, we first specify the
set of $Q$ primes dividing a site. When this set has bounded size,
its possible values give a manageable expansion into residue events.

For an offset $a$, its \emph{active $Q$-set} is
\[
             Q_a=\{p\in Q:R_p+a\equiv0\pmod p\}.
\]
Write $m_Q=\lfloor400\log L\rfloor$.  There are at most
\begin{equation}\label{q:q-state-count}
 \sum_{u=0}^{m_Q}\binom{|Q|}{u}
       \le(m_Q+1)(|Q|+1)^{m_Q}
       =\exp(O(L\log L))
\end{equation}
possible active sets with $|Q_a|\le m_Q$.  Specifying one exactly
is a conjunction testing every prime in $Q$ as present or absent.
On this event $w(n+a)=5^{|Q_a|}$, with $n$ denoting the random
origin as in the auxiliary model.  Every prescribed function or predicate
of this active set is then a scalar or a fixed decision.  In
particular, any sum over squarefree divisors of the product of the
active primes can be evaluated at this stage; there is no restriction
on how many such divisors are used in defining that scalar.

The event $|Q_a|>m_Q$ is also a small circuit: take the disjunction
over $(m_Q+1)$-element subsets of $Q$ of the conjunction asserting
that all their primes divide the site.  Define $\omega_P(n+a)=\#\{p\in P:R_p+a\equiv0\pmod p\}$.
The same construction applies
to the condition $\omega_P(n+a)>6WJ$, since
$6WJ\le\delta\log L$.  These circuits have size
$\exp(O(L\log L))$ and depth two.

The restriction $|Q_a|\le m_Q$ must precede an arbitrary predicate
of the active set. It reduces that predicate to a decision on each of
the states counted in \eqref{q:q-state-count}. We will apply this
representation to the padding cutoffs and weight denominators when
they arise. Its simplest consequence is the following weight bound.

\begin{corollary}[Truncated weight average]\label{q:truncated-weight}
For every offset $a$ and every interval $I$ allowed in
\Cref{q:finite-law},
\begin{equation}\label{q:truncated-weight-average}
 \E_I\bigl[w(n+a)\1_{\omega_Q(n+a)\le400\log L}\bigr]
                    \le S+e^{-L^9}\le2S
\end{equation}
for sufficiently large $L$.  The analogous sum over $M$ fixed offsets
is at most $2MS$.
\end{corollary}

\begin{proof}
Expand by exact active sets of size at most $m_Q$, with scalar
coefficient $5^{|Q_a|}$.  The total coefficient sum is
$\exp(O(L\log L))$, so \Cref{q:finite-expansion} applies.
The product-law expectation is at most $\E w=S$, by independence
of the prime coordinates.  Finally $S\ge1$.  The bound is uniform
in the offset, and summing it proves the last assertion.
\end{proof}

The comparison has now supplied the required passage to long integer
intervals for bounded residue tests and their explicit scalar
expansions.  Moments of the unrestricted weight will be used only
inside the auxiliary product law; the integer averages use the
truncated weight in \Cref{q:truncated-weight}.
